\subsection{代数的定义}

定义 1. 设 A 为一个交换环。A 上的代数（或 A-代数，或当不致混淆时简称代数）是一个集合 E，其结构由给出以下内容来定义：

\begin{enumerate}
\def\labelenumi{(\arabic{enumi})}
\item
  E 上的一个 A-模结构；
\item
  一个由 \(\mathbf{E} \times \mathbf{E}\) 到 \(\mathbf{E}\) 的 A-双线性映射（II，§ 3，第 5 号）。
\end{enumerate}

定义中出现的由 \(\mathbf{E} \times \mathbf{E}\) 到 \(\mathbf{E}\) 的 A-双线性映射称为代数 \(\mathbf{E}\) 的乘法；它通常记作 \((x, y) \mapsto x.y\) ，或简记作 \((x, y) \mapsto xy\) 。

设 \((\alpha_{i})_{i\in I}\) 与 \((\beta_{j})_{j\in J}\) 为 A 的两个有限支集的元素族（I，§ 2，第 1 号）。于是，对 E 的元素的所有族 \((x_{i})_{i\in I}\) 与 \((y_{j})_{j\in J}\) ，一般分配律公式（I，§ 3，第 4 号）

\[
\left(\sum_ {i \in I} \alpha_ {i} x _ {i}\right) \left(\sum_ {j \in J} \beta_ {j} y _ {j}\right) = \sum_ {(i, j) \in I \times J} \left(\alpha_ {i} \beta_ {j}\right) \left(x _ {i} y _ {j}\right)\tag{1}
\]

成立；特别地

\[
(\alpha x) y = x (\alpha y) = \alpha (x y) \quad \text { for } \alpha \in \mathrm{A}, x \in \mathrm{E} \text { and } y \in \mathrm{E}.\tag{2}
\]

双线性映射 \((x,y)\mapsto yx\) （由 E × E 到 E）与 E 上的 A-模结构在 E 上定义了一个 A-代数结构，称为与给定的代数结构相反的结构。带此新结构的集合 E 称为代数 E 的反代数；它常记作 \(E^{0}\) 。若 A-代数 E 与其反代数相同，则称 A-代数 E 为交换的，换言之，若 E 中的乘法是交换的。E 到 \(E^{0}\) 上的一个同构也称为代数 E 的反自同构。

当代数 E 中的乘法为结合的时，E 称为结合 A-代数。当 E 中的乘法具有单位元（它必然唯一（I，§ 2，第 1 号））时，此元素称为 E 的单位元，而 E 称为含单位元代数。

例子。(1) 每个交换环 A 都可以被视为一个（结合且交换的）A-代数。

\begin{enumerate}
\def\labelenumi{(\arabic{enumi})}
\setcounter{enumi}{1}
\item
  设 E 是一个伪环（I，§8，第 1 号）。E 上的乘法以及 E 上唯一的 Z-模结构在 E 上定义了一个结合的 Z-代数结构。
\item
  设 \(\mathbf{F}\) 为一个集合，\(\mathbf{A}\) 为一个交换环。集合 \(\mathbf{A}^{\mathbf{F}}\) （由 \(\mathbf{F}\) 到 \(\mathbf{A}\) 的所有映射组成），连同乘积环结构（I，§ 8，第 10 号）和乘积
\end{enumerate}

A-模结构（II，§ 1，第 5 号），是一个结合且交换的 A-代数。

\begin{enumerate}
\def\labelenumi{(\arabic{enumi})}
\setcounter{enumi}{3}
\tightlist
\item
  设 E 是一个 A-代数；内法则 \((x,y)\mapsto xy+yx\) 与 \((x,y)\mapsto xy-yx\) （连同 E 上的 A-模结构）在 E 上定义了两个 A-代数结构，它们一般不是结合的；第一个法则
\end{enumerate}

\[
(x, y) \mapsto x y + y x
\]

总是交换的。

定义 2. 给定交换环 A 上的两个代数 E, E'，E 到 E' 的一个同态是指一个映射 \(f: \mathbf{E} \to \mathbf{E}'\) ，使得

\begin{enumerate}
\def\labelenumi{(\arabic{enumi})}
\item
  \(f\) 是一个 A-模同态；
\item
  \(f(xy) = f(x)f(y)\) 对所有 \(x \in \mathbf{E}\) 和 \(y \in \mathbf{E}\) 成立。
\end{enumerate}

显然，两个 A-代数同态的复合仍是一个 A-代数同态。每个双射的代数同态都是一个同构。因此，A-代数同态可以取作 A-代数结构这一种类的态射（《集合论》，IV，2，第 1 号）。在下文中，我们总假定已经作出了这组态射的选取。若 E, E' 是两个 A-代数，用 \(\operatorname{Hom}_{\mathbf{A}-\operatorname{alg}}(\mathbf{E}, \mathbf{E}')\) 表示 E 到 E' 的 A-代数同态的集合。

设 E, E' 是两个均含单位元的代数。E 到 E' 的、把 E 的单位元映到 E' 的单位元的同态称为保单位同态（或保单位代数态射）。

